% TOP_PROBLEM: {"id": 2544, "title": "Kourovka Notebook Problem 21.35", "category_id": 17, "category": {"id": 17, "name": "group_theory", "display_name": "Group Theory", "description": "Problems about groups, group actions, representations, and related algebraic structures.", "slug": "group-theory", "order_index": 17, "created_at": "2026-07-31T15:26:25.671Z"}, "status": "open", "problem_number": "KOU-21.35", "set_id": 10, "statusQualification": "Primary v47 marks an unverified AI solution (#), expressly not a confirmed solution according to the editors; question retained."}
% FIRST_POSTED: 2026-10-01
% ENTRY_KIND: statement_only
% UnsolvedMath ID 2544; source catalogue code KOU-21.35
% !TeX program = lualatex
% Definition and sources only; this file is not an LLM solution attempt.
\documentclass[11pt,a4paper]{article}
\usepackage[margin=25mm]{geometry}
\usepackage{fontspec}
\setmainfont{lmroman10-regular.otf}[BoldFont=lmroman10-bold.otf,
 ItalicFont=lmroman10-italic.otf,BoldItalicFont=lmroman10-bolditalic.otf]
\usepackage{amsmath,amssymb,mathtools}
\usepackage{unicode-math}
\setmathfont{latinmodern-math.otf}
\AtBeginDocument{\let\setminus\smallsetminus}
\usepackage{xurl,hyperref}
\hypersetup{colorlinks=true,urlcolor=blue}
\setcounter{secnumdepth}{0}
\setlength{\parindent}{0pt}
\setlength{\parskip}{5pt}
\setlength{\emergencystretch}{3em}
\begin{document}
\section*{Kourovka Notebook Problem 21.35}\label{2544}
{\small UnsolvedMath ID 2544 · KOU-21.35 · Group Theory · Kourovka Notebook - New Problems, Issue 21}
\subsection{Definitions and mathematical statement}
Use the commutator convention $[u,v]=u^{-1}v^{-1}uv$.
A multilinear commutator word is obtained recursively from distinct
variables by repeatedly taking commutators of words involving disjoint
sets of variables. Thus every variable occurs once in its commutator
expression; different bracket arrangements are allowed.
For a word $w=w(x_1,\ldots,x_r)$ and a group $G$, a $w$-value is an
element $w(g_1,\ldots,g_r)$ with $g_i\in G$, and the verbal subgroup
$w(G)$ is generated by all such values.

Fix a finite group $G$, a prime $p$, and such a word $w$. For an element
$g$, let $o(g)$ be its order, and call this a $p'$-order if $p\nmid o(g)$.
Assume that for every pair of $w$-values $x,y$ in $G$,
\[
 p\nmid o(x),\quad p\mid o(y)\quad\Longrightarrow\quad p\mid o(xy).
\]
Must $w(G)$ be $p$-nilpotent? A finite group $K$ is $p$-nilpotent here
if it has a normal subgroup $N$ of order prime to $p$ with $K/N$ a
$p$-group; equivalently it has a normal Hall $p'$-subgroup.

The order hypothesis is imposed on individual word values, not on all
elements of the generated subgroup. The conclusion nevertheless concerns
the entire verbal subgroup. Both $x$ and $y$ are chosen from the same
word-value set, and their product need not itself be a word value.
Multilinearity of $w$ remains an essential hypothesis of the question.

The editors mark this entry with an unverified AI solution, which their
preface expressly distinguishes from a confirmed solution. The mathematical
question is therefore retained without endorsing that proposed argument.
\subsection{Short English statement}
Does the specified product-order condition on multilinear commutator values force their verbal subgroup to have a normal Hall $p'$-subgroup?
\subsection{Sources}
\begin{itemize}
\item[S1] ULAM.AI and the UnsolvedMath Contributors, supplied problems.json, record 2544 (KOU-21.35), Kourovka Notebook - New Problems, Issue 21. Catalogue identity and source transcription; CC BY 4.0.
\url{https://huggingface.co/datasets/ulamai/UnsolvedMath}.
\item[S2] The Kourovka Notebook, 21st edition, new problems, Problem 21.35. Expanded formulation of the supplied catalogue statement.
\url{https://arxiv.org/pdf/1401.0300v47}.
\end{itemize}
\end{document}
